% Drop-in replacement for the supplied draft's Simulation section.
% Required packages: amsmath, amssymb, graphicx, booktabs, array, multirow.
% Copy assets/ and tables/ alongside this fragment, or adjust graphicspath.
\providecommand{\simroot}{.}
\providecommand{\E}{\mathbb E}
\providecommand{\TV}{\operatorname{TV}}
\section{Numerical Results}
\label{sec:simulation}

We evaluate the framework through three complementary experiments. First,
image-classification comparisons examine deployment accuracy and bounded
risk under analog fading and additive noise. Second, controlled feature-erasure
experiments isolate the effect of a finite training channel bank and the
choice of training channel law. Finally, we evaluate the channel-shift,
channel-complexity, and data-complexity contributions at frozen learned
posteriors and compute a numerical population-risk bound. ERM and L2 are
the main competing baselines. Supplementary results record additional
internal ablations and sensitivity checks.

\subsection{Experimental Setup}
\label{sec:sim_setup}

\subsubsection{Data and split networks}
The image experiments use MNIST~\cite{mnist_dataset} and
Fashion-MNIST~\cite{fashion_mnist}. Both contain ten classes of
$28\times28$ grayscale images. For MNIST, we use 55,000 training and 5,000
validation examples from the official training partition. For Fashion-MNIST,
we use a subset of 10,000 training examples and a disjoint validation set of
5,000 examples. Both use the official 10,000-example test partition. Pixel
intensities are mapped to $[0,1]$, without augmentation or fitted
normalization. The common image classifier has dimensions $784\to128\to10$.
Its first affine layer and ReLU form the transmitter; its final affine
layer forms the receiver, with transmission between the two.

We additionally evaluate CIFAR-10~\cite{cifar10}, which contains ten classes
of $32\times32$ RGB images. Its official training partition is split into
45,000 training and 5,000 validation examples; the official test partition
contains 10,000 examples. Pixels are mapped to $[-1,1]$ by the fixed transform
$x/127.5-1$. The transmitter uses three $3\times3$ convolution/ReLU blocks
with 16, 32 and 64 output channels, max pooling after the first two blocks,
adaptive $4\times4$ pooling, and an affine/ReLU projection to 128 transmitted
features. A ten-class affine receiver follows the channel. All methods
share this 156,074-weight architecture, without batch normalization,
dropout or augmentation. This extends both the dataset and transmitter
architecture; it is not an isolated architecture ablation.

For controlled experiments, we sample an i.i.d. noisy two-moons population:
the binary label is Bernoulli with parameter $1/2$, the arc angle is uniform
on $[0,\pi]$, and independent Gaussian coordinate noise has standard
deviation $0.3$. These experiments use $2\to d\to2$ split networks. The
finite-channel study fixes $d=8$, $n=512$, a 512-example validation set, and
an independent 10,000-example test set. The training-law study uses $d=16$,
800 training examples, and 200 validation examples. Table~\ref{tab:sim_setup}
summarizes the main settings.

\begin{table*}[t]
\centering
\caption{Settings of the Main Experiments}
\label{tab:sim_setup}
\small
\setlength{\tabcolsep}{5pt}
\begin{tabular}{@{}lrrrrrrl@{}}
\toprule
Experiment & $n$ & Validation & Test & $d$ & $m$ & Epochs & Seeds \\
\midrule
MNIST, analog & 55,000 & 5,000 & 10,000 & 128 & 100 & 100 & 3 \\
Fashion-MNIST, analog & 10,000 & 5,000 & 10,000 & 128 & 100 & 25 & 2 \\
CIFAR-10, analog & 45,000 & 5,000 & 10,000 & 128 & 100 & 100 & 3 \\
Two moons, finite bank & 512 & 512 & 10,000 & 8 & $5$--$2000$ & 150 & 3 \\
Two moons, training-law mismatch & 800 & 200 & 10,000 & 16 & 100 & 150 & 3 \\
\bottomrule
\end{tabular}
\end{table*}

\subsubsection{Channel laws and finite observations}
For analog transmission, the recovered representation is
$\widehat f=Hf+N$, where $H$ is diagonal with independent
$H_{kk}\sim\mathcal N(\mu_m,\sigma_m^2)$ and
$N_k\sim\mathcal N(0,\sigma_b^2)$. The image comparisons use
$(\mu_m,\sigma_m,\sigma_b)=(0.5,1,1)$. These are fixed distortion-law
parameters: feature power is not normalized, so the comparisons are not at
a common physical SNR or transmit-power constraint. Neither CSI nor
equalization is supplied to the receiver.

For feature erasures, $\widehat f_k=B_kf_k$, where the coordinates
$B_k\sim\mathrm{Bernoulli}(1-p)$ are independent. The receiver observes
zero-filled features without a separate erasure-indicator input. This is
an erasure abstraction for continuous features, rather than a complete
bit-level digital transmission system. A channel-bank entry is one full
gain/noise vector or one full erasure mask. The bank is sampled separately
from the data and reused throughout training. Each training example in a
minibatch receives an independent uniform bank index, giving an unbiased
stochastic estimate of the empirical risk on $S\times U$.

\subsubsection{Bounded loss and optimization}
Let $p_w(y\mid x,c)$ be the confidence assigned to the correct class. With
$\varepsilon_\ell=0.01$, we use the common smooth bounded logarithmic loss
\begin{equation}
\label{equ:sim_bounded_loss}
\ell_{\varepsilon_\ell}(w,c,z)=
\frac{-\log[\varepsilon_\ell+(1-\varepsilon_\ell)p_w(y\mid x,c)]}
{\log(1/\varepsilon_\ell)}.
\end{equation}
It lies in $[0,1]$ uniformly over weights, channels, and data. Hence the
standing sub-Gaussian assumptions hold with
$\sigma^2=\sigma_{\rm tr}^2=1/4$ by Hoeffding's lemma. For perfect training,
the sharper channel variance proxy is zero.

The proposed learner uses one diagonal Gaussian posterior and the fixed,
data-independent prior $Q_W=\mathcal N(0,I)$. Its weight means start from
fan-in-scaled uniform initialization with zero biases; posterior standard
deviations initialize at $0.003$ or $0.01$ in the grayscale experiments,
and at $0.001$ or $0.003$ in the CIFAR-10 search.
We optimize with Adam and a constant learning rate, drawing one weight
sample per minibatch. Grayscale-image batches contain 256 examples;
CIFAR-10 and toy batches contain 128. The T1 objective uses the separate penalties in
Theorem~\ref{thm:separate_bound_US}; T2 uses the joint penalty in
Theorem~\ref{thm:joint_bound_US}. Unless a joint confidence allocation is
specified below, the nominal failure budget is $0.05$, split equally
between T1's two terms. ERM and L2 use the same loss
\eqref{equ:sim_bounded_loss}; L2 adds an explicit squared-weight penalty.

\subsubsection{Selection and evaluation}
The explored configurations are selected by mean validation accuracy over
the listed initialization/channel-bank seeds, breaking ties by lower
cross entropy. Checkpoint selection uses two fixed validation random-number
streams across epochs; test evaluation uses fresh streams. These are
exploratory numerical comparisons, and all listed seeds are retained in
the main tables. The data split is fixed across these seeds. Thus the
reported standard deviations measure initialization/channel-bank
variation conditional on that split. Search ranges and selected settings
are provided in the supplement.

Test accuracy is the Gibbs-expected accuracy: a weight draw is held fixed
for one complete test pass, with a fresh channel for each example, and
the resulting accuracy is averaged over ten passes. ERM and L2 are
deterministic, so only their channels vary. Bounded loss is evaluated
using the same posterior convention; ordinary cross entropy is a separate
metric. Perfect training means identity transmission in the training
forward pass, followed by testing through the deployment channel. For
Proposed Perfect, the objective retains only the data PAC-Bayes penalty.
Checkpoint selection for these references uses deployment-channel
validation; the perfect-training label refers to the training forward
channel.

\subsection{Deployment Performance Under Analog Fading}
\label{sec:sim_performance}
\input{\simroot/tables/accuracy_and_loss}

Table~\ref{tab:sim_accuracy_risk} reports mean test accuracy and normalized
bounded loss. On MNIST, T1 attains $76.63\%$, compared with $74.60\%$ for
matched ERM and $70.12\%$ for matched L2, giving gains of $2.03$ and $6.52$
percentage points. T2 reaches $76.43\%$ and slightly lowers the bounded
loss relative to T1. On Fashion-MNIST, T1 and T2 attain $63.24\%$ and
$63.21\%$, respectively, versus $61.57\%$ for ERM and $61.49\%$ for L2.
Both proposed variants outperform these evaluated baselines in every
reported seed of the two image comparisons. Their small differences in
accuracy and bounded loss do not imply a universal ordering between T1
and T2.

The perfect-training references also expose the importance of training
with channel observations. On MNIST, Proposed Perfect attains $57.94\%$
under deployment, whereas matched T1 attains $76.63\%$; the corresponding
Fashion-MNIST accuracies are $52.44\%$ and $63.24\%$. Removing the
finite-channel penalty during perfect training therefore does not remove
the deployment mismatch. These results concern the tested configurations;
they do not assert optimality over all baseline hyperparameters.

\subsubsection{CIFAR-10 with a convolutional transmitter}
\input{\simroot/tables/cifar_comparison}
Table~\ref{tab:sim_cifar} extends the comparison to the shared convolutional
architecture. The 34-run preliminary screen selects one setting per method
by aggregate validation accuracy. We then train these seven fixed settings
for 100 epochs on 45,000 examples using three new initialization/channel-bank
seeds; all three are retained. T1 achieves $55.03\%$ mean test accuracy,
versus $53.67\%$ for L2 and $46.16\%$ for ERM, gains of $1.36$ and $8.87$
percentage points. It exceeds both baselines in every reported seed and
also has the lowest mean bounded risk, $0.3395$, compared with $0.3526$
for L2 and $0.4205$ for ERM. T2 attains $53.26\%$, slightly below L2;
the benefit therefore does not extend uniformly to both proposed variants.
Mean transmitted feature power is $22.73$ per coordinate for T1, versus
$12.59$ for L2 and $6.71$ for ERM. These remain fixed-channel-law
comparisons with unconstrained feature power.

Perfect-trained ERM, L2 and Proposed attain approximately $34\%$ under
the same deployment law, reinforcing the training/deployment distinction.
The full-data comparison uses a fixed 100-epoch budget; several selected
checkpoints occur near its end, so these results do not establish fully
converged or globally optimal baseline accuracy. The preliminary screen
and search settings are recorded in the supplement. The high-dimensional
CNN bound at $m=100$ is vacuous; the numerical certificate below is
computed separately on the controlled small network.

\subsection{Generalization From a Finite Channel Bank}
\label{sec:sim_channel_bank}
For a frozen posterior, we evaluate bounded risk on the reused bank,
$\widehat R_{S,U}$, on the same data with fresh training-law channels,
$\widehat R_{S,\mathrm{tr}}$, and on held-out data with fresh training-law
channels, $\widehat R_{\mathrm{te},\mathrm{tr}}$. The signed channel gap is
\begin{equation}
\label{equ:sim_channel_gap}
\widehat G_{\rm ch}=\widehat R_{S,\mathrm{tr}}-\widehat R_{S,U}.
\end{equation}
This evaluates unseen channels without simultaneously changing the data.

\begin{figure*}[t]
\centering
\includegraphics[width=0.88\textwidth]{\simroot/assets/fig_channel_bank.pdf}
\caption{Effect of the reusable channel-bank size under matched feature
erasures: $p_{\rm tr}=p_{\rm ch}=0.5$, $n=512$, and $d=8$. Curves show
means over three initialization/channel-bank seeds; shading shows one
sample standard deviation and points retain individual seeds. The data,
number of optimizer updates, and deployment law are fixed across $m$.
Perfect-trained references are reused across bank sizes. The right panel
isolates the observed channel generalization gap in normalized-loss units.}
\label{fig:sim_channel_bank}
\end{figure*}

Figure~\ref{fig:sim_channel_bank} varies
$m\in\{5,20,100,500,2000\}$ while keeping the training budget at 600
minibatch updates. For T1, the observed channel gap decreases from
approximately $0.0103$ at $m=5$ to $0.0016$ at $m=20$, and is close to
zero from $m=100$ onward. This illustrates the additional generalization
effect induced by a finite channel bank even when the training and
deployment distributions match. The proposed variants nearly overlap in
this controlled example; the experiment supports the finite-channel
mechanism rather than a strict ranking between their optimizers. Residual
fluctuations at larger $m$ are comparable to the finite-evaluation scale.

\subsection{Choice of the Training Channel Distribution}
\label{sec:sim_mismatch}
\begin{figure*}[t]
\centering
\includegraphics[width=0.88\textwidth]{\simroot/assets/fig_training_mismatch.pdf}
\caption{Fixed deployment feature-erasure probability $p_{\rm ch}=0.5$
and varying training probability, with $n=800$, $d=16$, and $m=100$.
Perfect training is the point $p_{\rm tr}=0$; its data-only PAC objective
is common to the proposed perfect-channel specialization. Means and one
sample standard deviation use all three listed seeds. The vertical line
marks matched training.}
\label{fig:sim_training_mismatch}
\end{figure*}

Figure~\ref{fig:sim_training_mismatch} fixes deployment and considers
perfect training together with training erasure probabilities $0.1$,
$0.3$, $0.5$, $0.7$, and $0.9$. T1 achieves mean accuracies
of $83.13\%$, $84.42\%$, $84.90\%$, $84.84\%$, and $84.27\%$ at the
five nonzero probabilities. The matched point also minimizes cross entropy
in this experiment. Proposed Perfect achieves $81.87\%$, about three
percentage points below matched T1. Thus a perfect training forward
channel does not provide the best deployment accuracy in this setting.
Mildly and severely mismatched training laws yield different results even
though the deployment law remains unchanged. The theory guarantees that
matching removes the shift penalty; it does not require matched training
to maximize accuracy for every model and finite dataset.

\subsection{Bound Decomposition and Numerical Risk Guarantees}
\label{sec:sim_bounds}
The held-out deployment estimate is denoted by
$\widehat R_{\mathrm{te},\mathrm{ch}}$. We decompose its difference from
the reused-bank empirical estimate into
\begin{align}
\widehat G_{\rm data}
 &=\widehat R_{\mathrm{te},\mathrm{tr}}-\widehat R_{S,\mathrm{tr}},\nonumber\\
\widehat G_{\rm shift}
 &=\widehat R_{\mathrm{te},\mathrm{ch}}-\widehat R_{\mathrm{te},\mathrm{tr}},\nonumber\\
\widehat R_{\mathrm{te},\mathrm{ch}}-\widehat R_{S,U}
 &=\widehat G_{\rm ch}+\widehat G_{\rm data}+\widehat G_{\rm shift}.
\label{equ:sim_signed_decomposition}
\end{align}
Every term uses the same frozen posterior and paired weight draws. These
signed estimates are distinct from the nonnegative upper-bound terms.

For the bounded loss, channel shift also satisfies
$|\E_P g-\E_Q g|\leq\TV(P,Q)$ for $0\leq g\leq1$. For independent
feature erasures, total variation is evaluated exactly by grouping masks
by their number of erased coordinates. Perfect training therefore gives
the finite shift penalty $1-(1-p_{\rm ch})^d$, despite infinite channel
KL when $p_{\rm ch}>0$. This bounded-loss specialization does not require
an estimated channel-Lipschitz constant.

\begin{figure*}[t]
\centering
\includegraphics[width=0.88\textwidth]{\simroot/assets/fig_bound_decomposition.pdf}
\caption{Unweighted bound terms and signed observed gaps for perfect,
matched, and mismatched training, at the same deployment law
$p_{\rm ch}=0.5$, $n=800$, $d=16$, and $m=100$. Each column evaluates
the frozen posterior learned under its stated training law. Stacked
theoretical terms are those of T1; markers evaluate T2 at the same
posterior. The per-run nominal failure budget is $0.05$. Shift uses the
bounded-loss total-variation specialization. These illustrative gap
estimates do not include the additional MC correction used in
Table~\ref{tab:sim_certificate}.}
\label{fig:sim_bound_decomposition}
\end{figure*}

Figure~\ref{fig:sim_bound_decomposition} separates the mechanisms. Perfect
training has no channel complexity but retains deployment mismatch.
Matched training removes the shift penalty while retaining finite-channel
and data complexity. A stochastic mismatched training law generally
retains all three terms. The signed measured gaps need not be positive.
The displayed gap bounds exceed one and are vacuous on the normalized-loss
scale; their purpose here is to distinguish the three contributions.

\input{\simroot/tables/certificate_table}
We additionally compute a corrected population-risk bound for a chosen
T1-trained posterior with $n=512$, $m=2000$, and $d=8$, using the same
posterior for both theorem evaluations. A fresh 100-draw empirical-risk
estimate is upper-bounded by Hoeffding's inequality with failure budget
$0.01$. To cover selection among the seven declared matched-erasure
$(n,m)$ settings and three bank seeds, we assign $0.04/(2\times21)$ to each theorem/setup
event. The union of these events and the single fresh MC event has failure
probability at most $0.05$.

The corrected population-risk upper bounds are $0.8760$ for T1 and
$0.7178$ for T2, including MC error. Both are below one. The joint bound
is tighter at the same posterior; it does not change the predictor or
imply an accuracy ordering. These bounds concern the normalized Gibbs
loss in \eqref{equ:sim_bounded_loss}, rather than classification error or unbounded CE.
